% !TeX program = pdflatex
% Довідник для підготовки до екзамену з курсу
% "Чисельне моделювання динаміки систем".
% Компіляція: pdflatex dovidnyk.tex (двічі, для змісту).
\documentclass[12pt,a4paper]{article}

\usepackage[T2A]{fontenc}
\usepackage[utf8]{inputenc}
\usepackage[english,ukrainian]{babel}
\usepackage{lmodern}

\usepackage{amsmath,amssymb,amsthm}
\usepackage{mathtools}
\usepackage[a4paper,margin=2.3cm]{geometry}
\usepackage{enumitem}
\usepackage{microtype}
\usepackage{booktabs}
\usepackage{xcolor}
\usepackage{titlesec}
\usepackage[hidelinks,colorlinks=true,linkcolor=blue!55!black,urlcolor=teal!70!black]{hyperref}
\usepackage{fancyhdr}

% --- Оформлення заголовків ---
\definecolor{accent}{RGB}{20,80,140}
\titleformat{\section}
  {\normalfont\Large\bfseries\color{accent}}{\thesection.}{0.6em}{}
\titleformat{\subsection}
  {\normalfont\large\bfseries\color{accent!85!black}}{\thesubsection.}{0.5em}{}

\pagestyle{fancy}
\fancyhf{}
\renewcommand{\headrulewidth}{0.3pt}
\fancyhead[L]{\small\itshape Чисельне моделювання динаміки систем}
\fancyhead[R]{\small\thepage}

% --- Середовища для теорем/означень ---
\theoremstyle{definition}
\newtheorem{ozn}{Означення}[section]
\newtheorem{vlast}{Властивість}[section]
\theoremstyle{plain}
\newtheorem{teorema}{Теорема}[section]
\newtheorem*{teorema*}{Теорема}

% Блок «коротко» — резюме відповіді
\newenvironment{korotko}
  {\par\medskip\noindent\begingroup\color{accent!75!black}\rule{\linewidth}{0.4pt}\par
   \noindent\textbf{Коротко.}\itshape\,}
  {\par\noindent\rule{\linewidth}{0.4pt}\endgroup\medskip}

% --- Зручні позначення ---
\newcommand{\dd}[2]{\dfrac{\partial #1}{\partial #2}}
\newcommand{\ddt}[1]{\dfrac{\partial #1}{\partial t}}
\newcommand{\ddx}[1]{\dfrac{\partial #1}{\partial x}}
\newcommand{\ddxx}[1]{\dfrac{\partial^2 #1}{\partial x^2}}
\newcommand{\R}{\mathbb{R}}
\newcommand{\norm}[1]{\left\lVert #1 \right\rVert}
\newcommand{\abs}[1]{\left\lvert #1 \right\rvert}
\newcommand{\Real}{\operatorname{Re}}
\newcommand{\Imag}{\operatorname{Im}}
\DeclareMathOperator{\sgn}{sgn}
\DeclareMathOperator{\divv}{div}
\DeclareMathOperator{\grad}{grad}

\begin{document}

% ======================= ТИТУЛЬНА СТОРІНКА =======================
\begin{titlepage}
\centering
\vspace*{1.5cm}
{\scshape\large Київський національний університет імені Тараса Шевченка\par}
{\scshape факультет комп'ютерних наук та кібернетики\par}
\vspace{0.4cm}
{\itshape кафедра обчислювальної математики\par}
\vspace{2.2cm}
{\color{accent}\rule{\linewidth}{1.2pt}\par}
\vspace{0.6cm}
{\Huge\bfseries Чисельне моделювання\\[0.25em] динаміки систем\par}
\vspace{0.5cm}
{\LARGE Довідник для підготовки до екзамену\par}
\vspace{0.35cm}
{\large відповіді на екзаменаційні питання\par}
\vspace{0.4cm}
{\color{accent}\rule{\linewidth}{1.2pt}\par}
\vspace{2.2cm}
{\large
Спеціальність: прикладна математика\\[0.3em]
10 семестр\par}
\vfill
{\small Довідник укладено на основі курсу лекцій та рекомендованої літератури\\
(П.~Роуч; Р.~Рихтмайер, К.~Мортон; А.~А.~Самарський; О.~Ю.~Грищенко, С.~І.~Ляшко).\par}
\vspace{0.6cm}
{\small \today\par}
\end{titlepage}

% ======================= ЗМІСТ =======================
\tableofcontents
\newpage

\section*{Як користуватися довідником}
\addcontentsline{toc}{section}{Як користуватися довідником}
\noindent
Довідник містить структуровані відповіді на всі 11 екзаменаційних питань
курсу \emph{«Чисельне моделювання динаміки систем»}. Кожен розділ
відповідає одному питанню білета; на початку розділу подано блок
\textbf{«Коротко»} зі стислою суттю відповіді (зручно для повторення),
далі~--- докладний розбір з формулами, доведеннями та прикладами.
Наприкінці наведено список використаних джерел.

\medskip
\noindent
Нагадаємо структуру екзаменаційного білета: він містить
\emph{теоретичне питання} (одне з 11 наведених нижче) та
\emph{практичну задачу}. Найчастіше у практичній частині потрібно
дослідити задану різницеву схему \emph{на стійкість методом фон Неймана},
але можуть бути і задачі на \emph{порядок апроксимації},
\emph{дисипативність} та \emph{дисперсійність} схеми. Тому, окрім теорії,
у розділах 6--11 наведено робочі приклади дослідження схем, а в окремому
розділі~12 \textbf{«Практикум»}~--- уніфіковану методику та повністю
розв'язані типові практичні задачі. Їх варто опрацювати для практичної
частини.

\newpage

% =====================================================================
\section{Мета, задачі і проблеми чисельного моделювання. Рівняння гідродинаміки}
% =====================================================================

\begin{korotko}
Чисельне моделювання (обчислювальний експеримент)~--- третій метод
дослідження поряд із теорією та натурним експериментом. Базовою моделлю
гідродинаміки в'язкої нестисливої рідини є система рівнянь
Нав'є--Стокса разом із рівнянням нерозривності. Її записують у
\emph{недивергентній} (конвективній) та \emph{дивергентній}
(консервативній) формах; перехід до безрозмірних змінних виділяє єдиний
визначальний параметр~--- число Рейнольдса $\mathrm{Re}$.
\end{korotko}

\subsection{Мета, задачі і проблеми}
При дослідженні реальних фізичних процесів використовують три
взаємодоповнюючі підходи: \emph{теоретичний}, \emph{експериментальний}
та \emph{обчислювальний експеримент}. Чисельне моделювання бере на себе
частину функцій перших двох: воно дозволяє замінити високовартісні чи
неможливі натурні експерименти розрахунком на ЕОМ, прогнозувати
поведінку системи та, варіюючи параметри, виявляти нові
закономірності.

\smallskip\noindent\textbf{Доцільність.} Чисельне моделювання
застосовують, коли процеси важко або неможливо відтворити на макетах:
велика швидкоплинність, надмалі розміри, багатокомпонентність, надвисокі
чи наднизькі тиски й температури тощо.

\smallskip\noindent\textbf{Основні \emph{проблеми} (вимоги до методу).}
Система рівнянь повинна адекватно описувати процес і мати ефективні
\emph{обґрунтовані} чисельні методи. Обґрунтування методу полягає у
встановленні: існування розв'язку, оцінки похибки, обчислювальної
\emph{стійкості} алгоритму, його \emph{збіжності} до точного розв'язку,
а також \emph{консервативності}, \emph{дисипативності},
\emph{дисперсійності} тощо. Обчислювальна гідродинаміка~--- окрема
дисципліна зі своїми методами, труднощами та сферою застосування
(аерокосмічна галузь, автомобіле- і кораблебудування, біомедицина,
метеорологія, енергетика та ін.).

\subsection{Основні закони та рівняння}
Рух нестисливої ньютонівської в'язкої рідини в декартовій (ейлеровій)
системі координат описують \textbf{рівняння Нав'є--Стокса} (закон
збереження кількості руху, друге начало Ньютона до елемента об'єму) і
\textbf{рівняння нерозривності} (закон збереження маси):
\begin{align}
  \ddt{\mathbf{U}} + (\mathbf{U}\cdot\grad)\mathbf{U}
    &= -\frac{1}{\rho}\grad P + \nu\,\Delta\mathbf{U} + \mathbf{F}, \label{eq:ns-vec}\\
  \divv \mathbf{U} &= 0, \label{eq:cont}
\end{align}
де $\mathbf{U}$~--- вектор швидкості, $P$~--- тиск, $\rho$~--- густина,
$\nu$~--- кінематичний коефіцієнт в'язкості, $\mathbf{F}$~--- масові
сили. У двовимірному випадку $(u,v)$:
\begin{align}
  \ddt{u}+u\ddx{u}+v\dd{u}{y} &= -\frac1\rho\ddx{P}
     +\nu\Big(\ddxx{u}+\dd{^2 u}{y^2}\Big)+f_x,\\
  \ddt{v}+u\ddx{v}+v\dd{v}{y} &= -\frac1\rho\dd{P}{y}
     +\nu\Big(\ddxx{v}+\dd{^2 v}{y^2}\Big)+f_y,\\
  \ddx{u}+\dd{v}{y} &= 0 .
\end{align}

\paragraph{Змінні «функція току~--- вихор».} Вводячи вихор
$\omega=\dd{u}{y}-\ddx{v}$ і функцію току $\psi$ ($u=\dd{\psi}{y}$,
$v=-\ddx{\psi}$), із \eqref{eq:ns-vec} вилученням тиску дістають
\emph{рівняння переносу вихору} (параболічне за часом) та
\emph{рівняння Пуассона} (еліптичне):
\begin{equation}
  \ddt{\omega}+(\mathbf{U}\cdot\grad)\omega=\nu\,\Delta\omega,
  \qquad \Delta\psi=-\omega .
  \label{eq:vort}
\end{equation}

\subsection{Дивергентна і недивергентна форми}
\begin{itemize}[leftmargin=1.4em]
  \item \textbf{Недивергентна (конвективна)} форма містить конвективний
  член у вигляді $(\mathbf{U}\cdot\grad)\omega$.
  \item \textbf{Дивергентна (консервативна)} форма містить його у вигляді
  дивергенції потоку. За допомогою тотожності
  $\divv(\omega\mathbf{U})=\omega\,\divv\mathbf{U}+(\mathbf{U}\cdot\grad)\omega$
  та умови нерозривності $\divv\mathbf{U}=0$ рівняння \eqref{eq:vort}
  набуває вигляду
  \begin{equation}
     \ddt{\omega}+\divv(\omega\mathbf{U})=\nu\,\Delta\omega .
  \end{equation}
\end{itemize}
Дивергентна форма важлива тим, що різницеві схеми, побудовані на її
основі, природно зберігають інтегральні закони збереження
(консервативність, див. розд.~2 і 8).

\subsection{Обезрозмірювання рівнянь}
Нехай $L$~--- характерна довжина, $U_0$~--- характерна швидкість.
Уводимо безрозмірні величини
\[
  u^\ast=\frac{u}{U_0},\quad v^\ast=\frac{v}{U_0},\quad
  x^\ast=\frac{x}{L},\quad y^\ast=\frac{y}{L},\quad
  t^\ast=\frac{t}{L/U_0}.
\]
Тоді рівняння переносу вихору й Пуассона набувають вигляду (зірочки
опускаємо)
\begin{equation}
  \ddt{\omega}+\divv(\mathbf{V}\omega)=\frac{1}{\mathrm{Re}}\,\Delta\omega,
  \qquad \Delta\psi=-\omega,
\end{equation}
де визначальний безрозмірний параметр~--- \textbf{число Рейнольдса}
\begin{equation}
  \boxed{\;\mathrm{Re}=\dfrac{U_0 L}{\nu}\;}.
\end{equation}
Отже, за фіксованих граничних умов течія характеризується одним
параметром $\mathrm{Re}$. При $\mathrm{Re}\gg1$ переважає конвекція, при
$\mathrm{Re}\ll1$~--- дифузія; для останнього випадку зручнішим є
«дифузійний» масштаб часу $t^\ast=t/(L^2/\nu)$, що усуває виродження
рівняння при $\mathrm{Re}\to0$.

\paragraph{Одновимірні модельні рівняння.} Для дослідження схем
використовують модельні рівняння, що зберігають ключові риси
Нав'є--Стокса:
\begin{align}
  \text{лінійне (конвекція{+}дифузія):}\quad
    & \ddt{\varphi}+u\ddx{\varphi}=\alpha\ddxx{\varphi}, \label{eq:model}\\
  \text{рівняння Бюргерса:}\quad
    & \ddt{u}+u\ddx{u}=\alpha\ddxx{u}.
\end{align}
Тут $\alpha\sim1/\mathrm{Re}$~--- узагальнений коефіцієнт дифузії.
Рівняння Бюргерса зберігає нелінійність і є модельним для турбулентності
та ударних хвиль.

\newpage
% =====================================================================
\section{Методи апроксимації диференціальних операторів. Властивості схем}
% =====================================================================

\begin{korotko}
Різницеві аналоги похідних будують чотирма способами: розкладанням у
\emph{ряд Тейлора}, \emph{інтерполяційними многочленами},
\emph{інтегро-інтерполяційним методом} та \emph{методом контрольного
об'єму}. Усі вони можуть давати однаковий вираз, але залишають свободу
вибору. Ключові властивості одержаних схем~--- \emph{консервативність},
\emph{транспортивність}, \emph{дисипативність}.
\end{korotko}

\subsection{Метод рядів Тейлора}
Нехай $f$ достатньо гладка, крок сітки $\Delta x=h$. З розкладів
$f_{i\pm1}=f_i\pm h f_x+\tfrac{h^2}{2}f_{xx}\pm\tfrac{h^3}{6}f_{xxx}+\dots$
дістаємо різницеві апроксимації першої похідної:
\begin{align}
  \text{вперед:}\;& \delta_x f=\frac{f_{i+1}-f_i}{h}=f_x+O(h),\\
  \text{назад:}\;& \delta_x f=\frac{f_i-f_{i-1}}{h}=f_x+O(h),\\
  \text{центральна:}\;& \delta_x f=\frac{f_{i+1}-f_{i-1}}{2h}=f_x+O(h^2),
\end{align}
та другої похідної з другим порядком точності
\begin{equation}
  \delta_x^2 f=\frac{f_{i-1}-2f_i+f_{i+1}}{h^2}=f_{xx}+O(h^2).
\end{equation}
П'ятиточковий аналог рівняння Лапласа ($\beta=\Delta x/\Delta y$):
\[
  f_{i-1,j}-2(1+\beta^2)f_{i,j}+f_{i+1,j}
  +\beta^2(f_{i,j-1}+f_{i,j+1})=0 .
\]
Підставляючи центральні різниці у модельне рівняння \eqref{eq:model},
дістають, зокрема, схему \textbf{ВЧЦП} (FTCS, різниці вперед за часом,
центральні~--- по простору)
\begin{equation}
  \frac{\varphi_i^{n+1}-\varphi_i^n}{\Delta t}
  +u\,\frac{\varphi_{i+1}^n-\varphi_{i-1}^n}{2h}
  =\alpha\,\frac{\varphi_{i-1}^n-2\varphi_i^n+\varphi_{i+1}^n}{h^2},
  \label{eq:ftcs}
\end{equation}
яка має точність $O(\Delta t,h^2)$ і умовно стійка (див.~розд.~5--7).
Натомість схема з центральними різницями і за часом (рівн. типу
«чехарда» з самого початку) безумовно нестійка~--- ілюстрація того, що
точна апроксимація \emph{кожного} члена не гарантує придатної схеми.

\subsection{Поліноміальна (інтерполяційна) апроксимація}
Будують інтерполяційний многочлен за значеннями у вузлах і аналітично
диференціюють його. Для параболи $f=a+bx+cx^2$ за трьома точками
$i-1,i,i+1$ дістають у точці $i$ ті самі центральні формули другого
порядку. Поліноми вищого порядку чутливі до «шумів» у даних, тому в
обчислювальній гідродинаміці їх застосовують переважно для апроксимації
похідних поблизу границь.

\subsection{Інтегро-інтерполяційний (інтегральний) метод}
Рівняння у консервативній формі $\dd{\varphi}{t}=-\dd{(u\varphi)}{x}+\alpha\dd{^2\varphi}{x^2}$
інтегрують по комірці $[x-\tfrac h2,x+\tfrac h2]\times[t,t+\Delta t]$.
Точно інтегруючи по одній зі змінних і застосовуючи теорему про середнє
та формулу прямокутників для решти інтегралів, отримують різницеву
схему. Залежно від вибору точок усереднення можна одержати як схему
\eqref{eq:ftcs}, так і інші (зокрема безумовно нестійку). Метод
особливо зручний у непрямокутних системах координат.

\subsection{Метод контрольного об'єму}
Найбільш «фізичний» метод. Для контрольного об'єму (КО) з центром у
вузлі записують баланс шуканої величини $\Gamma=\bar\varphi\cdot V$:
\[
  \text{приріст }\Gamma\text{ у КО}=
  \text{конвективний потік крізь грані}+\text{дифузійний потік крізь грані}.
\]
Конвективний потік крізь грань $i+\tfrac12$ за консервативною формою
$\tfrac12(u_i\varphi_i+u_{i+1}\varphi_{i+1})$ \emph{точно} дорівнює
потоку, що втікає у сусідній КО,~--- звідси консервативність. Дифузійний
потік оцінюють за законом Фіка $q=-\alpha\,\partial\varphi/\partial x$.

\paragraph{Висновок.} Усі чотири методи можуть приводити до однакових
різницевих виразів, але в кожному є свобода вибору, тому
\emph{диференціальне рівняння не визначає різницевий аналог однозначно}.
Існує багато правдоподібних, але непрацездатних схем.

\subsection{Основні властивості різницевих схем}

\begin{ozn}[Консервативність]
Різницева схема \emph{консервативна}, якщо вона тотожно (на сітковій
множині) забезпечує виконання інтегральних законів збереження, що
справджуються для вихідного диференціального рівняння.
\end{ozn}
Для модельного рівняння при $\alpha=0$ ($\varphi_t+(u\varphi)_x=0$)
сума по області
$\sum_i \tfrac{\varphi_i^{n+1}-\varphi_i^{n}}{\Delta t}h$
зводиться до різниці потоків лише на границях (телескопічне
скорочення)~--- це сітковий аналог формули Остроградського--Гаусса.
Для \emph{недивергентної} форми $\varphi_t+u\varphi_x=0$ члени
взаємно не скорочуються (крім випадку $u=\mathrm{const}$), і виникають
штучні джерела~--- схема неконсервативна.

\begin{ozn}[Транспортивність]
Схема \emph{транспортивна}, якщо збурення, накладене у вузлі,
переноситься нею лише в напрямі вектора швидкості (за потоком).
\end{ozn}
\textbf{Приклад.} Центрально-різницева схема для
$\varphi_t+k\varphi_x=0$ переносить ізольоване збурення $\varepsilon$
\emph{в обидва боки}:
$\varepsilon_{i+1}^{n+1}=-\tfrac12 C\varepsilon$,
$\varepsilon_{i-1}^{n+1}=+\tfrac12 C\varepsilon$ (де $C=k\Delta t/h$)~---
тому вона \emph{не транспортивна}. Схема з різницями проти потоку
($k>0$: $\varphi_t+k\,(\varphi_i-\varphi_{i-1})/h=0$) переносить
збурення лише за потоком ($\varepsilon_{i+1}^{n+1}=C\varepsilon$,
$\varepsilon_{i-1}^{n+1}=0$)~--- транспортивна. \emph{Усі неявні схеми
не транспортивні}, бо реалізуються розв'язанням системи по всіх вузлах
шару (усереднення по всьому шару).

\paragraph{Схеми з донорними комірками.} Різновид схем проти потоку, де
знак середньої швидкості на грані визначає, з якого вузла брати значення
$\varphi$; вони транспортивні й консервативні та мають дещо вищий
порядок апроксимації просторової похідної.

\begin{ozn}[Дисипативність]
\emph{Дисипативність}~--- здатність схеми згладжувати (демпфувати)
гармоніки розв'язку; пов'язана з похідними \emph{парного} порядку у
виразі похибки апроксимації (штучна в'язкість). Детально~--- у розд.~11.
\end{ozn}

\newpage
% =====================================================================
\section{Точний та узагальнений розв'язки. Розв'язуючі оператори. Теорема Лакса}
% =====================================================================

\begin{korotko}
Стан системи~--- точка банахового простору $B$; еволюція~--- дія
\emph{розв'язуючого оператора} $E(t)$. Задача коректна за Адамаром, якщо
область визначення $E_0(t)$ щільна в $B$ і сім'я $\{E_0(t)\}$ рівномірно
обмежена. Для різницевої схеми $u^{n+1}=C(\Delta t)u^n$
\textbf{теорема Лакса} (про еквівалентність): для коректної задачі та
узгодженої апроксимації \emph{стійкість необхідна і достатня для
збіжності}.
\end{korotko}

\subsection{Точний та узагальнений розв'язки диференціальної задачі}
Нехай $B$~--- нормований банахів простір; функції від $x$ при фіксованому
$t$ ототожнюємо з точками $u\in B$. Розглядаємо задачу Коші
\begin{equation}
  \ddt{u}=Au,\qquad u(0)=u_0,\quad 0\le t\le T,
  \label{eq:cauchy}
\end{equation}
$A$~--- лінійний оператор.
\begin{ozn}
\emph{Точний розв'язок}~--- сім'я $u(t)$, для якої $u(t)$ належить
області визначення $A$, $u(0)=u_0$ і
$\big\|(u(t+\Delta t)-u(t))/\Delta t-Au(t)\big\|\to0$ рівномірно за $t$
при $\Delta t\to0$.
\end{ozn}
Позначимо через $D$ множину тих $u_0$, для яких точний розв'язок існує і
єдиний. Оператор $E_0(t):D\to B$, $u(t)=E_0(t)u_0$,~--- \textbf{точний
розв'язуючий оператор}.

\begin{ozn}[Коректність за Адамаром]
Задача коректна, якщо: \;1) $D$ щільна в $B$; \;2) сім'я $E_0(t)$
рівномірно обмежена: $\exists K>0$: $\norm{E_0(t)}\le K$, $0\le t\le T$.
\end{ozn}
Умова~1 означає, що будь-яку початкову умову можна як завгодно точно
наблизити «хорошими» (з $D$); умова~2~--- неперервну залежність
розв'язку від початкових даних.

\paragraph{Узагальнений розв'язуючий оператор.} За теоремою про
розширення обмежений лінійний оператор $E_0(t)$ зі щільною областю
визначення має єдине обмежене розширення $E(t)$ на весь $B$ з тією самою
сталою $K$. Тоді $u(t)=E(t)u_0$~--- \textbf{узагальнений розв'язок} для
будь-якого $u_0\in B$. Для нестаціонарного $A$ маємо
$E(t_2,t_0)=E(t_2,t_1)E(t_1,t_0)$ (\emph{напівгрупова властивість}).
Для неоднорідної задачі $u_t=Au+g(t)$ справджується формула
\begin{equation}
  u(t)=E(t)u_0+\int_0^t E(t-s)g(s)\,ds .
\end{equation}

\subsection{Різницеві задачі. Узгодженість, збіжність, стійкість}
Однорідне різницеве рівняння $B_1 u^{n+1}=B_0 u^n$ за існування
обмеженого $B_1^{-1}$ зводять до
\begin{equation}
  u^{n+1}=C(\Delta t)\,u^n,\qquad C(\Delta t)=B_1^{-1}B_0 .
  \label{eq:rec}
\end{equation}
\begin{ozn}[Узгодженість (апроксимація)]
Сім'я $C(\Delta t)$ \emph{узгоджено апроксимує} \eqref{eq:cauchy}, якщо
для щільної множини точних розв'язків
\[
  \Big\|\frac{C(\Delta t)-I}{\Delta t}\,u(t)-Au(t)\Big\|\to0,
  \quad \Delta t\to0 .
\]
\end{ozn}
\begin{ozn}[Стійкість]
Апроксимація \emph{стійка}, якщо для деякого $\tau>0$ множина операторів
$\{C(\Delta t)^n:\ 0\le n\Delta t\le T,\ 0<\Delta t\le\tau\}$ рівномірно
обмежена.
\end{ozn}
\begin{ozn}[Збіжність]
Для кожного $u_0\in B$ і кожної послідовності $\Delta_j t\to0$ з
$n_j\Delta_j t\to t$:\quad
$\big\|C(\Delta_j t)^{n_j}u_0-E(t)u_0\big\|\to0$.
\end{ozn}

\begin{teorema*}[Лакса про еквівалентність]
Нехай задача \eqref{eq:cauchy} коректно поставлена, а різницева
апроксимація \eqref{eq:rec} задовольняє умову узгодженості. Тоді
\emph{стійкість необхідна і достатня для збіжності}.
\end{teorema*}
\begin{proof}[Схема доведення]
\textbf{Необхідність.} Якби збіжна схема була нестійкою, знайшлися б
послідовності $\Delta_j t\to0$, $n_j$ з $\norm{C(\Delta_j t)^{n_j}u_0}\to\infty$,
причому $n_j\Delta_j t\to t\in[0,T]$; але збіжність вимагає
$C(\Delta_j t)^{n_j}u_0\to E(t)u_0$~--- суперечність. Отже множина
операторів рівномірно обмежена (стійкість).

\noindent\textbf{Достатність.} Для точного розв'язку $u(t)=E(t)u_0$ з
щільного класу оцінюють похибку
\[
  \xi_j=\Big(C(\Delta_j t)^{n_j}-E(n_j\Delta_j t)\Big)u_0
   =\sum_{k=0}^{n_j-1}C(\Delta_j t)^{k}\big(C(\Delta_j t)-E(\Delta_j t)\big)
     E\big((n_j-1-k)\Delta_j t\big)u_0 .
\]
За узгодженістю кожний доданок є $o(\Delta_j t)$ рівномірно, тож
$\norm{\xi_j}\le K_2\,n_j\Delta_j t\le K_2 T\,\varepsilon\to0$. Заміна
$E(n_j\Delta_j t)\to E(t)$ і щільність класу дають збіжність для всіх
$u_0\in B$.
\end{proof}

\paragraph{Приклад (принцип максимуму).} Неявна схема для рівняння
дифузії
$\dfrac{u_j^n-u_j^{n-1}}{\Delta t}=\dfrac{u_{j-1}^n-2u_j^n+u_{j+1}^n}{h^2}$
задовольняє принцип максимуму: внутрішній максимум $u_j^n$ не перевищує
максимуму початкових і граничних значень. Звідси $u_j^n$ обмежені при
будь-якій сітці~--- схема стійка.

\newpage
% =====================================================================
\section{Поняття стійкості та збіжності різницевих схем}
% =====================================================================

\begin{korotko}
\emph{Збіжність}~--- прямування різницевого розв'язку до точного при
$h,\Delta t\to0$ (через апроксимацію всього розв'язку, а не окремих
похідних). \emph{Стійкість}~--- рівномірна обмеженість степенів
розв'язуючого оператора $C(\Delta t)^n$ (обмеженість зростання
збурень/похибок округлення). За теоремою Лакса для коректної узгодженої
задачі ці поняття еквівалентні. Розрізняють \emph{динамічну} і
\emph{статичну} нестійкості.
\end{korotko}

\subsection{Апроксимація, збіжність, стійкість}
\begin{ozn}[Апроксимація/збіжність (Лакс, Рихтмайер)]
Збіжна різницева схема~--- та, що дає розв'язок, який збігається до
розв'язку диференціального рівняння при прямуванні кроків сітки до нуля
(збіжність розуміється для \emph{всього} розв'язку).
\end{ozn}
\begin{ozn}[Стійкість]
Схема стійка, якщо встановлено межу, до якої може зростати будь-яка
компонента початкових даних у процесі обчислень; формально~--- рівномірна
обмеженість $\{C(\Delta t)^n\}$ (розд.~3).
\end{ozn}
\textbf{Зв'язок.} Теорема Лакса: \emph{коректність + узгодженість}
$\Rightarrow$ (стійкість $\Leftrightarrow$ збіжність). Тому на практиці
достатньо перевірити апроксимацію (узгодженість) та стійкість.

\subsection{Похибка апроксимації та порядок точності}
Підставляючи точний розв'язок у різницеве рівняння, дістають
\emph{похибку апроксимації} $\psi^n$. Схема має \textbf{порядок точності}
$(p,q)$, якщо $\psi^n=O(\Delta t^{\,p}+h^{q})$. Наприклад, схема FTCS
\eqref{eq:ftcs} має порядок $O(\Delta t,h^2)$.

\subsection{Фізична картина нестійкості}
Розглянемо стаціонарний розв'язок $\hat\varphi^n$ з накладеним збуренням
$\varepsilon$. Для схеми FTCS приріст збурення складається з
конвективної та дифузійної частин, пропорційний $\Delta t$:
\begin{itemize}[leftmargin=1.4em]
  \item \textbf{Динамічна нестійкість}~--- осциляції зростаючої
  амплітуди через надто великий крок $\Delta t$ (дифузійний член
  «перекорегує» збурення). Усувається зменшенням $\Delta t<\Delta t_{кр}$.
  \item \textbf{Статична нестійкість}~--- монотонне зростання похибки
  через центральні різниці у конвективному члені; \emph{не} усувається
  зменшенням $\Delta t$, лише зміною схеми (напр., різниці проти потоку).
\end{itemize}
Якщо присутні обидва члени~--- вони взаємодіють, і обмеження на
$\Delta t$ залежить від співвідношення конвекції й дифузії, тобто від
сіткового числа Рейнольдса $\mathrm{Re}_C=uh/\alpha$.

\subsection{Методи дослідження стійкості (огляд)}
\begin{enumerate}[leftmargin=1.6em]
  \item \emph{Метод дискретних збурень} (розд.~6)~--- прямий, прослідковує
  загасання локального збурення.
  \item \emph{Метод фон Неймана (Фур'є)} (розд.~7)~--- через множник/матрицю
  переходу; для задач зі сталими коефіцієнтами.
  \item \emph{Метод Хьорта / диференціального наближення} (розд.~7, 11)~---
  через еквівалентне диференціальне рівняння.
  \item \emph{Енергетичний метод} (розд.~9)~--- для змінних коефіцієнтів;
  дає достатні умови.
  \item \emph{Принцип заморожених коефіцієнтів, ознака Бабенка--Гельфанда}
  (розд.~10).
\end{enumerate}

\newpage
% =====================================================================
\section{Багатошарові схеми та багатокрокові алгоритми}
% =====================================================================

\begin{korotko}
$(q{+}1)$-шарова схема зводиться до системи двошарових уведенням нових
змінних і \emph{матричного} розв'язуючого оператора $C$. Дослідження
(узгодженість, стійкість, збіжність) проводять у розширеному просторі
вектор-стовпчиків. Особлива проблема~--- задання початкових умов:
багатошарова схема потребує значень на кількох шарах, тому початкові дані
доводять окремою (двошаровою) процедурою.
\end{korotko}

\subsection{Зведення до одношарової (двошарової) схеми}
Задачу $u_t=Au$ апроксимуємо $(q{+}1)$-шаровою схемою
\begin{equation}
  B_q u^{n+q}+B_{q-1}u^{n+q-1}+\dots+B_0 u^n=0,
  \label{eq:multilayer}
\end{equation}
яку (за оборотності $B_q$) переписують як
$u^{n+q}=C_{q-1}u^{n+q-1}+\dots+C_0 u^n$, $C_i=-B_q^{-1}B_i$. Уводячи
$v^n=u^{n+1},\ w^n=v^{n+1},\dots$, дістаємо систему двошарових рівнянь,
яку записують у векторному вигляді
\begin{equation}
  \Phi^{n+1}=C\,\Phi^n,\qquad
  \Phi^n=\big(u^{n+q-1},\dots,u^{n}\big)^{\!T},
\end{equation}
де \textbf{матриця переходу} (супровідна матриця)
\begin{equation}
  C=\begin{pmatrix}
    C_{q-1} & C_{q-2} & \cdots & C_1 & C_0\\
    E & 0 & \cdots & 0 & 0\\
    0 & E & \cdots & 0 & 0\\
    \vdots & & \ddots & & \vdots\\
    0 & 0 & \cdots & E & 0
  \end{pmatrix}.
\end{equation}
Норму у розширеному просторі вводять, наприклад, як
$\norm{(u,v,\dots)}=(\norm{u}^2+\norm{v}^2+\dots)^{1/2}$.

\subsection{Постановка початкових умов}
Диференціальна задача має одну початкову умову в $t_0$, а різницева
\eqref{eq:multilayer}~--- значення на $q$ шарах. Якщо їх обчислити точно
($\Phi^0=S u_0$, де $S=\mathrm{diag}(E((q{-}1)\Delta t),\dots,E,\dots)$),
то наближення $\Phi^n=C^n S u_0$ апроксимує $E(n\Delta t)Su_0$. На
практиці $\Phi^0$ отримують \emph{апроксимацією} (часто допоміжним
двошаровим методом з меншим кроком, щоб не втратити порядок точності).
Для збіжності потрібно, щоб \emph{незалежно} $u^n\to u(t)$ і
$\Phi^0\to u(0)$ при $\Delta t\to0$.

\subsection{Властивості та критерії}
\begin{ozn}
Багатошарова схема \emph{стійка}, якщо сім'я $\{C^n(\Delta t):0\le n\Delta t\le T,\,0<\Delta t\le\tau\}$
рівномірно обмежена.
\end{ozn}
Справджуються твердження (Рихтмайер--Мортон):
\begin{enumerate}[leftmargin=1.6em]
  \item якщо схема стійка і похибка апроксимації $\to0$ при $\Delta t\to0$,
  то виконується умова узгодженості;
  \item для коректної задачі та узгодженої багатошарової апроксимації
  \emph{стійкість необхідна й достатня для збіжності} (аналог теореми
  Лакса).
\end{enumerate}
\textbf{Спектральний критерій.} Для стійкості необхідно, щоб
спектральний радіус матриці переходу
$\rho(G)=\max_j\abs{\lambda_j}\le1+O(\Delta t)$ (умова фон Неймана для
багатошарових схем; приклад~--- тришарова схема, розд.~7).

\paragraph{Недолік багатошарових (зокрема тришарових) схем.} Потреба у
кількох початкових шарах робить різницеве рівняння рівнянням вищого
порядку й може спричиняти перевизначеність; це компенсують допоміжним
розрахунком стартових шарів.

\newpage
% =====================================================================
\section{Метод дискретних збурень}
% =====================================================================

\begin{korotko}
Прямий метод: у вузол вводять дискретне збурення $\varepsilon$ і
прослідковують його еволюцію крок за кроком. Схема стійка, якщо збурення
\emph{загасають} ($\abs{\varepsilon^{n+1}/\varepsilon^{n}}\le1$). Метод
наочний, придатний і для перевірки \emph{транспортивності}, але стає
громіздким на дальших шарах і дає лише необхідні умови.
\end{korotko}

\subsection{Ідея методу}
У знайдений (нехай стаціонарний, $\varphi_i^n\equiv0$) розв'язок вводять
у вузлі $i$ збурення $\varepsilon$ і за різницевою схемою обчислюють
збурення на наступних шарах та в сусідніх вузлах; вимога стійкості~---
їх загасання.

\subsection{Рівняння дифузії, схема FTCS}
Для $\varphi_t=\alpha\varphi_{xx}$ зі схемою
$\varphi_i^{n+1}=\varphi_i^n+d(\varphi_{i-1}^n-2\varphi_i^n+\varphi_{i+1}^n)$,
\;$d=\alpha\Delta t/h^2$ (\emph{дифузійне число}), збурення $\varepsilon$
у вузлі $i$ дає
\[
  \varepsilon_i^{n+1}=(1-2d)\varepsilon,\qquad
  \varepsilon_{i\pm1}^{n+1}=d\,\varepsilon .
\]
Вимога $\abs{\varepsilon_i^{n+1}}\le\abs{\varepsilon}$ дає
$-1\le1-2d\le1$, тобто $d\le1$. Розгляд наступних шарів і найгіршого
(осцилюючого) розподілу збурень $\varepsilon_i=(-1)^i\varepsilon$ дає
$\varepsilon_i^{n+1}=(1-4d)\varepsilon$ і остаточну умову
\begin{equation}
  \boxed{\,d=\frac{\alpha\Delta t}{h^2}\le\tfrac12\,}\quad\Longleftrightarrow\quad
  \Delta t\le\frac{h^2}{2\alpha}.
  \label{eq:diff-stab}
\end{equation}
(Вимога відсутності осциляцій $\varepsilon_i^{n+1}/\varepsilon\ge0$ дає
ту саму межу $d\le\tfrac12$.) Обмеження \eqref{eq:diff-stab} «дороге»:
здрібнення $h$ удвічі у $D$-вимірній задачі збільшує машинний час у
$2^{\,D}\cdot 4$ разів (для $D=1$~--- у $8$ разів).

\subsection{Конвекція{+}дифузія}
Для \eqref{eq:ftcs} (нехай $u<0$) аналіз у вузлах $i\pm1$ дає
\[
  \varepsilon_{i\mp1}^{n+1}=\Big(\pm\tfrac{C}{2}+d\Big)\varepsilon,
  \qquad C=\frac{u\Delta t}{h},\ \ d=\frac{\alpha\Delta t}{h^2},
\]
звідки крім \eqref{eq:diff-stab} дістають умову статичної стійкості
$\Delta t\le 2\alpha/u^2$, а умова відсутності осциляцій~--- обмеження на
сіткове число Рейнольдса
\begin{equation}
  \mathrm{Re}_C=\frac{\abs{u}h}{\alpha}\le2 .
\end{equation}
У комбінації це дає число Куранта $C=\abs{u}\Delta t/h\le1$ та
$\Delta t\le2\alpha/u^2$. (Метод фон Неймана дає дещо іншу комбіновану
умову $C^2\le2d\le1$.)

\paragraph{Завдання-зразок.} Для схеми проти потоку
$(\varphi_i^{n+1}-\varphi_i^n)/\Delta t=-u(\varphi_i^n-\varphi_{i-1}^n)/h$,
$u>0$, метод дискретних збурень дає
$\varepsilon_i^{n+1}=(1-C)\varepsilon$, $\varepsilon_{i+1}^{n+1}=C\varepsilon$,
$\varepsilon_{i-1}^{n+1}=0$ (транспортивність!), а умова стійкості~---
\emph{число Куранта} $C=u\Delta t/h\le1$.

\newpage
% =====================================================================
\section{Метод фон Неймана. Коефіцієнт і матриця переходу}
% =====================================================================

\begin{korotko}
Розв'язок шукають у вигляді ряду Фур'є; кожна гармоніка
$u^n_j=\sum_m A_m V^n(m)e^{Imjh}$ за лінійності схеми зі сталими
коефіцієнтами множиться на \textbf{множник переходу} $G(m)$ (для систем~---
\textbf{матрицю переходу} $G=H_1^{-1}H_0$). \textbf{Необхідна умова фон
Неймана}: $\rho(G)=\max_j\abs{\lambda_j}\le1+O(\Delta t)$. Для нормальної
$G$ (зокрема для скалярного $G$ двошарової схеми) умова \emph{необхідна й
достатня}.
\end{korotko}

\subsection{Подання Фур'є та умова стійкості}
Періодичну $u(x)$ розкладають у ряд Фур'є; рівність Парсеваля
$\int_0^L\abs{u}^2dx=\sum_{k}\abs{\hat u(k)}^2$ задає ізоморфізм
$B\leftrightarrow\tilde B$ зі збереженням норми. Різницеве рівняння
$B_1u^{n+1}=B_0u^n$ у Фур'є-образах:
\begin{equation}
  H_1\,\eta^{n+1}(k)=H_0\,\eta^{n}(k),\qquad
  H_j=\sum_{\nu\in N_j} B_\nu\,e^{I(k_1\nu_1\Delta x_1+\dots)},
\end{equation}
звідки $\eta^{n+1}(k)=G(k)\,\eta^n(k)$, $G=H_1^{-1}H_0$~--- матриця
переходу. \textbf{Умова стійкості}~--- рівномірна обмеженість
$\{G^n(k)\}$. Необхідна спектральна умова:
\begin{equation}
  \boxed{\;\rho(G)=\max_j\abs{\lambda_j(G)}\le1+M\Delta t\;}
\end{equation}
Оскільки $\rho(G)\le\norm{G}$, то достатньою є $\norm{G}\le1+M\Delta t$
(тоді $\norm{G^n}\le e^{MT}$). Для \emph{нормальної} матриці
($\rho(G)=\norm{G}_2$) умова фон Неймана необхідна й достатня; зокрема для
скалярного множника $p=1$ (двошарова схема, одна залежна змінна).

\begin{teorema*}[про збіжність ряду Фур'є]
Якщо $\norm{G(m)}<1$ і початкова функція розвивається в абсолютно збіжний
ряд Фур'є, то точний періодичний розв'язок різницевого рівняння зі
сталими коефіцієнтами подається збіжним рядом Фур'є
$u^n_j=\sum_m A_m V^n(m)e^{Imjh}$.
\end{teorema*}
\begin{proof}[Ідея]
Помноживши $V^{n+1}=GV^n$ на $A_m e^{Imjh}$ і підсумувавши, дістають
$\norm{u^{n+1}}\le q^n\sum_m\abs{A_m}$, $q=\max_m\abs{G(m)}<1$~---
мажоранта збіжності; ряд задовольняє рівняння та початкові/граничні
умови.
\end{proof}

\subsection{Явна схема для рівняння теплопровідності}
Для $u_t=\alpha u_{xx}$ явна схема
$\dfrac{u_i^{n+1}-u_i^n}{\Delta t}=\alpha\dfrac{u_{i-1}^n-2u_i^n+u_{i+1}^n}{h^2}$
підстановкою $u^n_j=V^n e^{Imjh}$ дає
\begin{equation}
  G(m)=V(m)=1+d\big(\cos mh-1\big)=1-2d\sin^2\tfrac{mh}{2},\quad d=\frac{\alpha\Delta t}{h^2}.
\end{equation}
Умова $\abs{G}\le1$ для всіх $m$:\;
$-1\le1-2d\sin^2\tfrac{mh}{2}\le1\Rightarrow d\le\tfrac12$, тобто
\begin{equation}
  \boxed{\;\Delta t\le\frac{h^2}{2\alpha}\;}\quad(\text{умовна стійкість}).
\end{equation}

\begin{teorema*}[стійкість явної схеми]
Явна схема для рівняння теплопровідності стійка тоді й лише тоді, коли
$d=\alpha\Delta t/h^2\le\tfrac12$.
\end{teorema*}

\subsection{Неявна схема для рівняння теплопровідності}
Для
$\dfrac{u_i^{n+1}-u_i^n}{\Delta t}=\alpha\dfrac{u_{i-1}^{n+1}-2u_i^{n+1}+u_{i+1}^{n+1}}{h^2}$
дістаємо
\begin{equation}
  G(m)=\frac{1}{1+2d\sin^2(mh/2)} .
\end{equation}
\begin{teorema*}[стійкість неявної схеми]
$\abs{G(m)}<1$ для всіх $m$ і будь-яких $\Delta t,h$~--- неявна схема
\emph{безумовно стійка}.
\end{teorema*}

\subsection{Рівняння конвективного переносу}
Для $u_t=-k u_x$ \emph{центральна} явна схема дає
$G=1+IC\sin mh$, $\abs{G}^2=1+C^2\sin^2 mh>1$~--- \emph{безумовно
нестійка}. \emph{Схема проти потоку} (перший порядок) дає множник з
$\abs{G}\le1$ за умови \textbf{Куранта} $C=\abs{k}\Delta t/h\le1$.

\paragraph{Тришарова схема (матриця переходу).} Для $u_t=-k u_x$
тришарова схема
$\dfrac{u_i^{n+1}-u_i^{n-1}}{2\Delta t}=-k\dfrac{u_{i+1}^n-u_{i-1}^n}{2h}$
(другого порядку) зводиться до
$V^{n+1}=aV^n+V^{n-1}$, $a=-2IC\sin\theta$, $\theta=mh$, і векторного
рівняння з \textbf{матрицею переходу}
\[
  G(m)=\begin{pmatrix} a & 1\\ 1 & 0\end{pmatrix},\qquad
  \lambda^2-a\lambda-1=0,\ \ \lambda=\tfrac{a\pm\sqrt{a^2+4}}{2}.
\]
Якщо $C^2\sin^2\theta>1$, корені чисто уявні з $\abs{\lambda}>1$
(нестійкість); якщо ж $C\le1$, то $\abs{\lambda_1}=\abs{\lambda_2}=1$~---
схема стійка при $C=1$ і дає при $C=1$ точний розв'язок.

\newpage
% =====================================================================
\section{Стійкість схем для дифузії та конвекції. Консервативність і транспортивність}
% =====================================================================

\begin{korotko}
Зведення воєдино: для рівняння \emph{дифузії} ($u_t=\alpha u_{xx}$) явна
схема стійка при $d\le\tfrac12$, неявна~--- безумовно. Для рівняння
\emph{конвекції} ($u_t=-ku_x$) центральна явна схема нестійка, схема
проти потоку стійка при числі Куранта $C\le1$. Багатовимірні випадки
дають жорсткіші обмеження. Окремо досліджують \emph{консервативність}
(сіткові закони збереження) та \emph{транспортивність} (перенос лише за
потоком).
\end{korotko}

\subsection{Зведена таблиця умов стійкості (метод фон Неймана)}
\begin{center}
\renewcommand{\arraystretch}{1.35}
\begin{tabular}{@{}p{5.2cm}p{4.4cm}p{4.6cm}@{}}
\toprule
\textbf{Рівняння / схема} & \textbf{Множник $G$} & \textbf{Умова стійкості}\\
\midrule
дифузія, явна (FTCS) & $1-2d\sin^2\tfrac{mh}{2}$ & $d=\alpha\Delta t/h^2\le\tfrac12$\\
дифузія, неявна & $\big(1+2d\sin^2\tfrac{mh}{2}\big)^{-1}$ & безумовно стійка\\
конвекція, центральна явна & $1+IC\sin mh$ & нестійка завжди\\
конвекція, проти потоку & $1-C(1-e^{-Imh})$ & $C=\abs{k}\Delta t/h\le1$\\
конвекція+дифузія (FTCS) & див. розд.~7 & $C^2\le2d\le1$\\
\bottomrule
\end{tabular}
\end{center}

\subsection{Багатовимірні задачі}
Для 2D рівняння переносу вихору (FTCS, $d_x=\alpha\Delta t/\Delta x^2$ та
ін.) множник переходу
\[
  G=1-2d_x(1-\cos\beta_x)-2d_y(1-\cos\beta_y)-I\big(C_x\sin\beta_x+C_y\sin\beta_y\big),
\]
$\beta_x=k_x\Delta x$, $\beta_y=k_y\Delta y$. Необхідні умови $\abs{G}\le1$
дають, зокрема, у випадку $d_x=d_y=d$:\; $d\le\tfrac14$ (удвічі жорсткіше,
ніж у 1D), а для 3D рівняння дифузії~--- $d\le\tfrac16$ (утричі жорсткіше).
Це підтверджує «дорожнечу» явних схем у багатовимірних задачах.

\subsection{Консервативність схем (детально)}
Консервативна схема тотожно зберігає на сітці інтегральний закон
збереження. Для $\varphi_t+(k\varphi)_x=0$ схема \emph{проти потоку}
($k>0$) консервативна: підсумовуючи по всіх вузлах, дістаємо
телескопічну суму
\[
  \Big(\sum_{i=1}^N \varphi_i^{n+1}h-\sum_{i=1}^N \varphi_i^{n}h\Big)\!\Big/\Delta t
  =-(k\varphi)^n_{N+1}+(k\varphi)^n_1,
\]
тобто зміна сумарної величини дорівнює різниці потоків на границях. Якщо
$k(x)$ змінює знак усередині області, поява членів типу
$(k\varphi)^n_{l+1}-(k\varphi)^n_l$ дає \emph{штучні джерела}~---
порушення консервативності. Його усувають донорною модифікацією схеми.

\subsection{Транспортивність схем (детально)}
Як показано у розд.~2 і 6, центральна схема не транспортивна (переносить
збурення в обидва боки), схема проти потоку та схема з донорними
комірками~--- транспортивні. Усі \emph{неявні} схеми не транспортивні. У
багатовимірних задачах забезпечення транспортивності значно складніше і
вимагає окремого дослідження для кожного випадку.

\newpage
% =====================================================================
\section{Стійкість зі змінними операторними коефіцієнтами. Енергетичний метод}
% =====================================================================

\begin{korotko}
Для рівнянь зі \emph{змінними коефіцієнтами} метод фон Неймана прямо не
застосовний. \textbf{Енергетичний метод} оперує спряженими різницевими
операторами та формулою підсумовування за частинами; домножуючи рівняння
на розв'язок (чи $u^{n+0,5}$) і підсумовуючи, дістають оцінку
$\norm{u^{n+1}}\le(1+K\Delta t)\norm{u^n}$~--- \emph{достатню} умову
стійкості. Метод дає більше інформації, ніж фон Неймана (зокрема
коректний результат для нестійкої центральної схеми).
\end{korotko}

\subsection{Огляд умов стійкості для змінних коефіцієнтів}
\begin{itemize}[leftmargin=1.4em]
  \item \textbf{Неперервна залежність порядку $p$}: задача коректна, якщо
  $\norm{u(\cdot,t)}\le \mathrm{const}\,\norm{u(\cdot,0)}_p$.
  \item \textbf{Сильна стійкість}: існує оператор $H(\Delta t)>0$ з
  $K_1^{-1}\norm{u}^2\le\norm{u}_H^2\le K_1\norm{u}^2$ та
  $\norm{u^{n+1}}_H\le(1+K_2\Delta t)\norm{u^n}_H$. Тоді
  $\norm{u^n}\le K_1 e^{K_2 T}\norm{u^0}_H$.
  \item \textbf{Локальна умова} (рівняння $u_t=a(x)u_{xx}$, $a(x)\ge\delta>0$):
  схема стійка, якщо $a$ ліпшицева, $\abs{G(x,\theta)}\le1$ і
  $G(x,0)=1$ (узгодженість).
  \item \textbf{Теорема Лакса--Ніренберга}: якщо матричні коефіцієнти
  $c_\nu(x)$ дійсні симетричні, з обмеженими другими похідними, і
  $\norm{\sum_\nu c_\nu(x)e^{I\nu\theta}}\le1$ всюди, то схема стійка:
  $\norm{C(\Delta t)}\le1+O(\Delta t)$.
\end{itemize}

\subsection{Спряжений оператор і підсумовування за частинами}
На просторі фінітних сіткових функцій зі скалярним добутком
$(u,v)=\sum_i u_i v_i h$ для оператора
$L_h u_i=k\dfrac{u_i-u_{i-1}}{h}$ формула підсумовування за частинами дає
спряжений
$L_h^\ast v_i=-k\dfrac{u_{i+1}-u_i}{h}$. Корисна тотожність:
\begin{equation}
  (L_h y,\,y)=\frac{h}{2k}\,\norm{L_h y}^2 .
\end{equation}

\subsection{Енергетична оцінка для одного рівняння}
Розглянемо $u_t=-k u_x$ ($k=\mathrm{const}>0$) зі схемою
$(u_i^{n+1}-u_i^n)/\Delta t=-L_h u_i^n$, де
$L_h u_i=k(u_i-u_{i-1})/h$. Тоді
$u^{n+0,5}=u^n-\tfrac{\Delta t}{2}L_h u^n$, і
\[
  \norm{u^{n+1}}^2-\norm{u^n}^2
   =-2\Delta t\,(L_h u^n,\,u^{n+0,5})
   =-2\Delta t\Big(\tfrac{h}{2k}-\tfrac{\Delta t}{2}\Big)\norm{L_h u^n}^2,
\]
де використано тотожність $(L_h u,u)=\tfrac{h}{2k}\norm{L_h u}^2$. Отже
\begin{equation}
  \norm{u^{n+1}}^2=\norm{u^n}^2-\Delta t\Big(\frac{h}{k}-\Delta t\Big)\norm{L_h u^n}^2 .
\end{equation}
Якщо права поправка недодатна, тобто
\begin{equation}
  \boxed{\;\Delta t\le \frac{h}{k}\;}\quad(\text{умова Куранта}),
\end{equation}
то $\norm{u^{n+1}}\le\norm{u^n}$~--- рівномірна стійкість за початковими
даними. \emph{Енергетичний метод дає достатню умову.}

\paragraph{Центральна схема через енергетичний метод.} Для оператора
$L_h^0 u_i=k\dfrac{u_{i+1}-u_{i-1}}{2h}$ маємо $(L_h^0 u,u)=0$, звідки
$\norm{u^{n+1}}^2\ge\norm{u^n}^2$ (норма не спадає), а оцінка зверху дає
$\norm{u^{n+1}}^2\le(1+k^2\Delta t/h^2)\norm{u^n}^2\le e^{0,5 d_0 T}\norm{u^0}^2$
з $d_0=k^2\Delta t/h^2$. Це «параболічна» (неприродна для гіперболічного
рівняння) умова, але вона коректно описує поведінку центральної схеми.

\subsection{Енергетичний метод для системи}
Для системи $u_t=k v_x$, $v_t=k u_x$ диференціальна енергія
$E(t)=\tfrac12\int(u^2+v^2)dx$ зберігається: $dE/dt=0$. Сімейству схем з
вагами $\sigma,\theta$ відповідає
\[
  J^{n+1}=J^n+\Delta t\,Q,\quad
  J=\tfrac12(\norm{u^n}^2+\norm{v^n}^2),\quad
  Q=-(\sigma-\tfrac12)\norm{L_{\Delta t}u^n}^2-(\theta-\tfrac12)\norm{L_{\Delta t}v^n}^2 .
\]
Якщо $Q\le0$~--- схема стійка в енергетичній нормі. Зокрема при
$\theta=\tfrac12$ і $\sigma\ge\tfrac12$ маємо $Q\le0$; при
$\sigma=\theta=\tfrac12$ $Q=0$~--- виконується сітковий закон збереження
енергії $\norm{u^{n+1}}^2+\norm{v^{n+1}}^2=\norm{u^0}^2+\norm{v^0}^2$;
при $\tfrac12<\sigma\le1$ енергія спадає (дисипація схеми, $O(\Delta t)$).

\newpage
% =====================================================================
\section{Принцип заморожених коефіцієнтів та ознака Бабенка--Гельфонда}
% =====================================================================

\begin{korotko}
\textbf{Принцип заморожених коефіцієнтів}: для стійкості лінійної
різницевої задачі зі змінними коефіцієнтами \emph{необхідно}, щоб у
кожній точці задача зі сталими «замороженими» коефіцієнтами задовольняла
спектральну ознаку фон Неймана; за робочу умову беруть найжорсткішу по
області. \textbf{Ознака Бабенка--Гельфонда} додатково враховує вплив
\emph{граничних умов}, досліджуючи три задачі: Коші та дві
напівобмежені (з лівою та з правою границею).
\end{korotko}

\subsection{Принцип заморожених коефіцієнтів}
У довільній точці $(\bar x,\bar t)$ фіксують значення коефіцієнтів і
вважають, що в усій області розв'язується рівняння зі сталими
коефіцієнтами; до нього застосовують ознаку фон Неймана.
\begin{teorema*}[принцип заморожених коефіцієнтів]
Для стійкості лінійної різницевої задачі зі змінними коефіцієнтами
необхідно, щоб кожна задача Коші зі сталими замороженими коефіцієнтами
задовольняла необхідну спектральну ознаку фон Неймана (\emph{локальна
стійкість}).
\end{teorema*}
\textbf{Приклад.} Для
$\dfrac{u_i^{n+1}-u_i^n}{\Delta t}=a(x,t)\dfrac{u_{i-1}^n-2u_i^n+u_{i+1}^n}{h^2}$
заморожування $a$ дає локальну умову
$\Delta t\le \dfrac{h^2}{2\,\bar a}$, а по всій області~---
\begin{equation}
  \Delta t\le\frac{h^2}{2\max_{x,t}a(x,t)} .
\end{equation}
Принцип придатний і евристично для нелінійних задач (напр.
$u_t=(1+u^2)u_{xx}$ дає прогноз кроку
$\Delta t_{p+1}\le h^2/\big(2\max_i(1+(u_i^n)^2)\big)$). Якщо локальну
необхідну умову порушено хоча б в одній точці~--- стійкості крайової
задачі очікувати не можна за жодних граничних умов. \emph{Обмеження:}
принцип не враховує вплив граничних умов; локальна стійкість може
залежати від них (приклади Крайса).

\subsection{Ознака Бабенка--Гельфонда}
Враховує граничні умови, вимагаючи дослідження \emph{трьох} різницевих
задач (з замороженими коефіцієнтами):
\begin{enumerate}[leftmargin=1.6em]
  \item \textbf{задача Коші} (без граничних умов)~--- збурення всередині
  області;
  \item \textbf{напівобмежена справа} ($i\ge0$): з урахуванням лише лівої
  граничної умови $l_1 u_0^{n+1}=0$, права границя віднесена в $+\infty$
  (для розв'язків з $u_i^n\to0$, $i\to+\infty$);
  \item \textbf{напівобмежена зліва} ($i\le M$): з урахуванням лише правої
  граничної умови $l_2 u_M^{n+1}=0$, ліва границя в $-\infty$.
\end{enumerate}
\begin{teorema*}[ознака Бабенка--Гельфонда]
Для стійкості початково-крайової задачі зі змінними коефіцієнтами
\emph{спектри операторів переходу всіх трьох задач повинні лежати в
замкненому одиничному крузі} $\abs{\lambda}\le1$.
\end{teorema*}

\subsection{Методика обчислення спектра}
Шукають розв'язки виду $u_i^n=\lambda^n u_i^0$. Для модельного рівняння
($a\equiv1$, $r=\Delta t/h^2$) підстановка $u_i^0=q^i$ дає
характеристичне рівняння
\begin{equation}
  q^2+\frac{-2r+1-\lambda}{r}\,q+1=0 .
\end{equation}
Для обмежених розв'язків $q=e^{I\theta}$ ($\abs{q}=1$):
\begin{equation}
  \lambda(\theta)=1-4r\sin^2\frac{\theta}{2},\qquad 0\le\theta\le2\pi,
\end{equation}
що заповнює відрізок $[1-4r,\,1]$~--- це спектр задачі Коші. Якщо
$\lambda\notin[1-4r,1]$, корені $q_1,q_2$ мають $\abs{q_1 q_2}=1$,
$\abs{q_1}<1<\abs{q_2}$; для напівобмеженої справа беруть спадний
$u_i=Cq_1^i$ і підставляють у граничну умову $l_1 u_0=0$, знаходячи
власні значення $\lambda$, що відповідають границі. Наприклад:
\begin{itemize}[leftmargin=1.4em]
  \item $l_1 u_0\equiv u_0=0$ або $u_1-u_0=0$: додаткових власних значень
  немає, спектр $=[1-4r,1]$;
  \item $l_1 u_0\equiv 2u_1-u_0=0$: дає $q_1=\tfrac12$ і власне значення
  $\lambda=1+r/2>1$~--- \emph{поза} одиничним кругом, тобто гранична умова
  спричиняє нестійкість.
\end{itemize}
Аналогічно досліджують третю задачу через $u_i=Cq_2^i$ та $l_2 u_M=0$.

\newpage
% =====================================================================
\section{Метод першого диференціального наближення. Схемна в'язкість}
% =====================================================================

\begin{korotko}
\textbf{Метод першого диференціального наближення (ПДН)}: усі сіткові
функції розкладають у ряд Тейлора, підставляють у схему й зберігають
члени з кроками у нульовому та першому степені~--- дістають
\emph{еквівалентне} диференціальне рівняння. Його аналіз (знак
коефіцієнтів) дає умови стійкості та виявляє \textbf{схемну (штучну)
в'язкість}. \emph{Гіперболічна} форма ПДН дає умову Куранта,
\emph{параболічна}~--- умову на дифузію. Похідні \emph{парного} порядку
у похибці $\Rightarrow$ \emph{дисипація}; \emph{непарного} $\Rightarrow$
\emph{дисперсія}.
\end{korotko}

\subsection{Ідея методу}
Оцінки апроксимації мають асимптотичний характер; реальні розрахунки
ведуться зі скінченними (хоч і малими) кроками, тому доцільно залишати
в апроксимуючому рівнянні принаймні перше наближення розкладу. Це
виявляє «приховані» в схемі властивості~--- штучну в'язкість, дисипацію,
дисперсію.

\subsection{Приклад~1: схема Лакса для $u_t=ku_x$}
Схема Лакса
$u_i^{n+1}=\tfrac12(u_{i+1}^n+u_{i-1}^n)+\tfrac{k\Delta t}{2h}(u_{i+1}^n-u_{i-1}^n)$.
Розкладаючи у ряд Тейлора, дістають \textbf{гіперболічну форму} ПДН
\begin{equation}
  \ddt{u}=k\ddx{u}+\frac{\Delta t}{2}\dd{^2u}{t^2}-\frac{h^2}{2\Delta t}\ddxx{u}+O(\dots),
\end{equation}
а виключивши $u_{tt}$ через $u_{tt}=k^2 u_{xx}$,~--- \textbf{параболічну
форму}
\begin{equation}
  \ddt{u}=k\ddx{u}+\nu_{\text{сх}}\ddxx{u},\qquad
  \boxed{\;\nu_{\text{сх}}=\frac12\Big(\frac{h^2}{\Delta t}-k^2\Delta t\Big)\;}.
\end{equation}
Тобто схема Лакса еквівалентна рівнянню \emph{дифузії} зі
\textbf{схемною в'язкістю} $\nu_{\text{сх}}$, відсутньою у вихідному
рівнянні. Коректність (невід'ємність дифузії) $\nu_{\text{сх}}\ge0$ дає
$h^2\ge k^2\Delta t^2$, тобто \textbf{умову стійкості} (Куранта)
\begin{equation}
  C=\frac{\abs{k}\Delta t}{h}\le1,
\end{equation}
причому при $C=1$ ($\nu_{\text{сх}}=0$) схема дає точний розв'язок.
Похибка схеми Лакса $O(\Delta t,h^2,h^2/\Delta t)$~--- схема апроксимує
рівняння лише якщо одночасно $\Delta t,h\to0$ \emph{і} $h^2/\Delta t\to0$.
Недолік: не транспортивна (переносить збурення проти потоку); перевага~---
простота й адаптивність до циліндричних/сферичних координат.

\subsection{Приклад~2: явна схема для $u_t+ku_x=\alpha u_{xx}$}
ПДН у \textbf{гіперболічній} формі дає область впливу з нахилом
характеристик $\pm\Delta t/(2\alpha)$, а умова Куранта (область впливу
схеми $\supseteq$ область впливу рівняння) дає
\begin{equation}
  \Delta t\le\frac{h^2}{2\alpha}.
\end{equation}
\textbf{Параболічна} форма (виключенням $u_{tt}$) дає
\begin{equation}
  \ddt{u}+k\ddx{u}=\tilde\alpha\,\ddxx{u},\quad \tilde\alpha=\alpha-\frac{k^2\Delta t}{2},
\end{equation}
звідки друга умова стійкості $\tilde\alpha\ge0$, тобто
$\Delta t\le 2\alpha/k^2$.

\subsection{Схемна (штучна) в'язкість, дисипація, дисперсія}
\begin{itemize}[leftmargin=1.4em]
  \item \textbf{Схемна в'язкість.} Схема проти потоку для $u_t+(cu)_x=0$
  еквівалентна рівнянню повного переносу з \emph{нефізичною} в'язкістю
  \begin{equation}
    \nu_c=\tfrac12 k h\,(1-C)\ge0 ,
  \end{equation}
  що згладжує розриви незалежно від причини градієнтів.
  \item \textbf{Дисипація}~--- наслідок похідних \emph{парного} порядку у
  похибці апроксимації; згладжує (демпфує) амплітуду, типова для схем
  \emph{першого} порядку; добра для стаціонарних задач, погана для
  розривів.
  \item \textbf{Дисперсія}~--- наслідок похідних \emph{непарного} порядку;
  спотворює фази хвиль (нефізичні осциляції перед/за фронтом), типова для
  схем \emph{другого} порядку; для гладких розв'язків.
  \item Спільну дію дисипації та дисперсії називають \emph{дифузією}
  розв'язку (розтяг крутих фронтів).
\end{itemize}

\subsection{Аналіз похибок через коефіцієнт переходу}
Записавши $G=\abs{G}e^{I\phi}$, порівнюють із точним
$G_e=e^{-Ik_m k\Delta t}$ ($\abs{G_e}=1$, $\phi_e=-k_m k\Delta t$):
\begin{itemize}[leftmargin=1.4em]
  \item \textbf{амплітудна (дисипативна) похибка} після $n$ кроків
  $\sim(1-\abs{G}^n)A_0$;
  \item \textbf{фазова (дисперсійна) похибка} $\phi/\phi_e$. Для схеми
  проти потоку при малих $\theta$:\;
  $\phi/\phi_e\approx1-\tfrac16(1-C)(1-2C)\theta^2$~--- при $0,5<C<1$
  випередження за фазою, при $C<0,5$~--- відставання.
\end{itemize}
Для \emph{неявного методу Ейлера} ПДН має вигляд
$u_t=ku_x+\tfrac12 k^2\Delta t\,u_{xx}$ (значна дисипація на середніх
хвильових числах і запізнення за фазою на великих).

\newpage
% =====================================================================
\section{Практикум: порядок апроксимації, дисипативність, дисперсійність}
% =====================================================================

\begin{korotko}
Для практичної задачі зручний єдиний алгоритм: підставити у схему
гармоніку $u_j^n=V^n e^{I\theta j}$ ($\theta=mh$) і знайти множник
переходу $G(\theta)$. Тоді: \;\textbf{стійкість}~--- $\abs{G}\le1$;
\;\textbf{порядок апроксимації}~--- з розкладу точного розв'язку у ряд
Тейлора (головний член похибки $O(\Delta t^{\,p}+h^{q})$);
\;\textbf{дисипативність}~--- з модуля $\abs{G(\theta)}$ (загасання
амплітуди; \emph{парні} похідні у похибці);
\;\textbf{дисперсійність}~--- з фази $\arg G(\theta)$ порівняно з точною
(спотворення фази; \emph{непарні} похідні у похибці).
\end{korotko}

\subsection{Загальний алгоритм дослідження схеми}
\begin{enumerate}[leftmargin=1.8em]
  \item \textbf{Множник переходу.} Підставити $u_j^n=V^n e^{I\theta j}$,
  $\theta=mh$, скоротити на $e^{I\theta j}$:\; $V^{n+1}=G(\theta)V^n$
  (для систем/багатошарових схем~--- матриця $G$).
  \item \textbf{Стійкість (фон Неймана):} $\abs{G(\theta)}\le1+O(\Delta t)$
  для всіх $\theta\in[-\pi,\pi]$ (для матриці~--- $\rho(G)\le1+O(\Delta t)$).
  \item \textbf{Порядок апроксимації:} підставити точний розв'язок
  $u(x,t)$, розкласти всі вузлові значення у ряд Тейлора в околі
  $(x_i,t_n)$; за потреби замінити похідні за $t$ на похідні за $x$ через
  саме рівняння. Головний член похибки $\psi=O(\Delta t^{\,p}+h^{q})$.
  \item \textbf{Дисипативність:} обчислити $\abs{G(\theta)}$. Якщо
  $\abs{G(\theta)}<1$ для $0<\abs{\theta}\le\pi$~--- схема \emph{дисипативна}
  (демпфує гармоніки); головний член похибки містить похідну
  \emph{парного} порядку (штучна в'язкість $\nu_{\text{сх}}u_{xx}$ тощо).
  \item \textbf{Дисперсійність:} обчислити фазу $\phi(\theta)=\arg G(\theta)$
  і порівняти з точною $\phi_e=-C\theta$. Відносна фазова похибка
  $\varepsilon_\phi=\phi/\phi_e$; головний член похибки містить похідну
  \emph{непарного} порядку ($u_{xxx}$ тощо).
\end{enumerate}

\begin{ozn}[Дисипативність за Крайсом]
Схему $u^{n+1}=C(\Delta t)u^n$ називають \emph{дисипативною порядку $2r$},
якщо існує $\delta>0$ таке, що
$\abs{\lambda(G)}\le1-\delta\,\abs{\theta}^{2r}$ для всіх
$\abs{\theta}\le\pi$. Тобто короткі хвилі ($\theta$ великі) загасають
сильніше за довгі.
\end{ozn}

\paragraph{Зв'язок з першим диференціальним наближенням (ПДН).} Якщо
головний член похибки апроксимації~--- похідна \emph{парного} порядку
($u_{xx},u_{xxxx},\dots$) з додатним коефіцієнтом, схема
\emph{дисипативна} (згасання амплітуди). Якщо \emph{непарного}
($u_{xxx},\dots$)~--- схема \emph{дисперсійна} (спотворення фази,
нефізичні осциляції). Схеми 1-го порядку зазвичай дисипативні, 2-го~---
дисперсійні (див. розд.~11).

\subsection{Задача A. Схема проти потоку для рівняння переносу}
\textbf{Умова.} Дослідити порядок апроксимації, дисипативність і
дисперсійність схеми проти потоку для $u_t+k u_x=0$, $k>0$:
\[
  \frac{u_i^{n+1}-u_i^n}{\Delta t}+k\,\frac{u_i^n-u_{i-1}^n}{h}=0,
  \qquad C=\frac{k\Delta t}{h}.
\]

\paragraph{1) Множник переходу.} $u_i^{n+1}=u_i^n-C(u_i^n-u_{i-1}^n)$,
підстановка $u_j^n=V^n e^{I\theta j}$ дає
\begin{equation}
  G(\theta)=1-C\big(1-e^{-I\theta}\big)
   =\big(1-C+C\cos\theta\big)-I\,C\sin\theta .
\end{equation}

\paragraph{2) Стійкість.} 
\begin{equation}
  \abs{G}^2=\big(1-C(1-\cos\theta)\big)^2+C^2\sin^2\theta
   =1-4C(1-C)\sin^2\tfrac{\theta}{2}.
\end{equation}
Для $0\le C\le1$ маємо $\abs{G}\le1$~--- схема \textbf{стійка при $C\le1$}
(умова Куранта).

\paragraph{3) Порядок апроксимації.} Підставляючи точний розв'язок і
розкладаючи у ряд Тейлора, дістаємо ПДН (розд.~11):
\begin{equation}
  u_t+k u_x=\underbrace{\frac{kh}{2}(1-C)}_{\nu_{\text{сх}}\ \ge 0}\,u_{xx}+O(h^2).
\end{equation}
Головний член похибки~--- $O(\Delta t,h)$, тобто схема має
\textbf{перший порядок} апроксимації.

\paragraph{4) Дисипативність.} Оскільки головний член ПДН містить
\emph{другу} (парну) похідну $u_{xx}$ зі схемною в'язкістю
$\nu_{\text{сх}}=\tfrac12 kh(1-C)\ge0$, а $\abs{G}<1$ для $0<C<1$ і
$\theta\ne0$,~--- схема \textbf{дисипативна} (демпфує амплітуду коротких
хвиль). При $C=1$ ($\nu_{\text{сх}}=0$, $\abs{G}=1$) дисипація зникає і
схема дає точний розв'язок. Амплітудна похибка за $n$ кроків
$\sim\big(1-\abs{G}^n\big)A_0$.

\paragraph{5) Дисперсійність.} Фаза
$\phi=\arg G=\arctan\dfrac{-C\sin\theta}{1-C(1-\cos\theta)}$, точна фаза
$\phi_e=-C\theta$. Розкладаючи для малих $\theta$:
\begin{equation}
  \varepsilon_\phi=\frac{\phi}{\phi_e}
   \approx 1-\frac16\,(1-C)(1-2C)\,\theta^2 .
\end{equation}
Отже схема \textbf{дисперсійна}: при $0{,}5<C<1$~--- випередження за фазою
($\varepsilon_\phi>1$), при $C<0{,}5$~--- відставання ($\varepsilon_\phi<1$).

\subsection{Задача B. Центральна (тришарова) схема «чехарда»}
\textbf{Умова.} Для $u_t+k u_x=0$ дослідити схему
\[
  \frac{u_i^{n+1}-u_i^{n-1}}{2\Delta t}+k\,\frac{u_{i+1}^n-u_{i-1}^n}{2h}=0 .
\]

\paragraph{Порядок.} Центральні різниці і за часом, і за простором дають
похибку $O(\Delta t^2,h^2)$~--- \textbf{другий порядок} апроксимації.

\paragraph{Множник переходу та стійкість.} $V^{n+1}=-2IC\sin\theta\,V^n+V^{n-1}$,
матриця переходу
$G=\left(\begin{smallmatrix}-2IC\sin\theta & 1\\ 1 & 0\end{smallmatrix}\right)$,
$\lambda^2+2IC\sin\theta\,\lambda-1=0$. При $C\le1$ обидва корені мають
$\abs{\lambda}=1$~--- схема \textbf{нейтрально стійка} ($C\le1$).

\paragraph{Дисипативність.} Оскільки $\abs{\lambda}=1$ для всіх $\theta$
(амплітуда не згасає), схема \textbf{бездисипативна}: головний член
похибки не містить похідних парного порядку.

\paragraph{Дисперсійність.} Головний член похибки~--- похідна
\emph{третього} (непарного) порядку $u_{xxx}$, тому схема
\textbf{сильно дисперсійна}: різні гармоніки рухаються з різними фазовими
швидкостями, що породжує нефізичні осциляції біля крутих фронтів
(характерно для схем 2-го порядку).

\subsection{Задача C. Схема Лакса}
\textbf{Умова.} Для $u_t+k u_x=0$:
$u_i^{n+1}=\tfrac12(u_{i+1}^n+u_{i-1}^n)-\tfrac{C}{2}(u_{i+1}^n-u_{i-1}^n)$.
\begin{itemize}[leftmargin=1.4em]
  \item \textbf{Множник переходу:} $G=\cos\theta-IC\sin\theta$,
  $\abs{G}^2=\cos^2\theta+C^2\sin^2\theta=1-(1-C^2)\sin^2\theta\le1$ при
  $C\le1$ (стійка).
  \item \textbf{Порядок:} $O(\Delta t,h^2,h^2/\Delta t)$~--- перший порядок
  (апроксимує лише при $\Delta t,h\to0$ і $h^2/\Delta t\to0$).
  \item \textbf{Дисипативність:} ПДН дає
  $\nu_{\text{сх}}=\tfrac12\big(h^2/\Delta t-k^2\Delta t\big)\ge0$ (парна
  похідна $u_{xx}$)~--- схема \textbf{дисипативна}; при $C=1$
  $\nu_{\text{сх}}=0$~--- точний розв'язок.
  \item \textbf{Дисперсійність:} наявна, але при першому порядку домінує
  дисипація.
\end{itemize}

\subsection{Зведена таблиця властивостей типових схем}
\begin{center}
\renewcommand{\arraystretch}{1.35}
\begin{tabular}{@{}p{4.6cm}p{2.7cm}p{2.5cm}p{2.5cm}@{}}
\toprule
\textbf{Схема (для $u_t+ku_x=0$)} & \textbf{Порядок} & \textbf{Дисипа\-тивність} & \textbf{Дисперсій\-ність}\\
\midrule
проти потоку (1-й пор.) & $O(\Delta t,h)$ & так & слабка\\
центральна явна (2-шар.) & --- & анти\-дисипативна & ---\\
«чехарда» (3-шарова) & $O(\Delta t^2,h^2)$ & ні & сильна\\
Лакса & $O(\Delta t,h^2,\tfrac{h^2}{\Delta t})$ & так & є\\
Лакса--Вендроффа & $O(\Delta t^2,h^2)$ & слабка & так\\
неявний Ейлер & $O(\Delta t,h)$ & сильна & запізнення фази\\
\bottomrule
\end{tabular}
\end{center}
\noindent\small Загальне правило: схеми \emph{1-го} порядку
переважно \emph{дисипативні} (згладжують розриви, добрі для стаціонарних
задач), схеми \emph{2-го} порядку~--- переважно \emph{дисперсійні}
(точніші на гладких розв'язках, але дають осциляції біля фронтів).

\newpage
% =====================================================================
\section*{Список використаних джерел}
\addcontentsline{toc}{section}{Список використаних джерел}
% =====================================================================
\begin{enumerate}[leftmargin=1.8em]
  \item \textbf{Курс лекцій} «Чисельне моделювання динаміки систем»
  (лекції 1--10), кафедра обчислювальної математики, КНУ ім.~Тараса
  Шевченка~--- основне джерело структури й змісту цього довідника.
  \item Роуч~П. \emph{Вычислительная гидродинамика}.~--- М.: Мир, 1980.~--- 616~с.
  \item Андерсон~Д., Танненхилл~Дж., Плетчер~Р. \emph{Вычислительная
  гидродинамика и теплообмен}. Т.~1--2.~--- М.: Мир, 1990.
  \item Самарський~А.~А. \emph{Теория разностных схем}.~--- М.: Наука, 1983.~--- 616~с.
  \item Грищенко~О.~Ю., Ляшко~С.~І. \emph{Методи Фур'є та першого
  диференціального наближення в теорії різницевих схем}.~--- К.: ВПЦ
  «Київський університет», 2005.~--- 84~с.\\
  {\small Онлайн (через веб-архів):
  \url{https://web.archive.org/web/2id_/http://www.unicyb.kiev.ua/Library/Methods_dif_sheems/zmist.htm}\\
  завантажені розділи див. у теці \texttt{sources/} (\texttt{methods\_*.htm}).}
  \item Рихтмайер~Р., Мортон~К. \emph{Разностные методы решения краевых
  задач}.~--- М.: Мир, 1972.~--- 418~с.
  \item Коллатц~Л. \emph{Функциональный анализ и вычислительная математика}.~--- М.: Мир.
  \item Годунов~С.~К., Рябенький~В.~С. \emph{Разностные схемы}.~--- М.: Наука, 1977.~--- 440~с.
  \item Марчук~Г.~И. \emph{Методы вычислительной математики}.~--- М.: Наука, 1989.~--- 384~с.
  \item Самарский~А.~А., Вабищевич~П.~Н. \emph{Численные методы решения
  задач конвекции--диффузии}.~--- М.: Эдиториал УРСС, 1999.~--- 248~с.
  \item Грищенко~О.~Ю., Ляшко~С.~І., Молодцов~О.~І. \emph{Чисельне
  моделювання процесів релаксаційної газової динаміки}.~--- К.: ІЗМН, 1997.~--- 224~с.\\
  {\small Онлайн (через веб-архів):
  \url{https://web.archive.org/web/20220419004248/http://www.unicyb.kiev.ua/Library/Calc_met_gas_dyn/Gr3.htm}\\
  завантажено у \texttt{sources/gas\_dyn\_Gr3.htm}.}
  \item Шокин~Ю.~И. \emph{Метод дифференциальных приближений}.~--- Новосибирск: Наука, 1979.~--- 219~с.
  \item Kreiss~H.~O. On difference approximations of the dissipative type
  for hyperbolic difference equations~// Comm. Pure Appl. Math.~--- 1964.~--- Vol.~17.~--- P.~355.
\end{enumerate}

\end{document}
